% =============================================================================
%  lec05.tex —— 第 5 讲 拟范畴的拟范畴（底本 PDF 第 39–49 页）
% =============================================================================
\lecture{5}{拟范畴的拟范畴}
\begin{center}\small\color{black!55}（底本 PDF 第 39–49 页）\end{center}

本讲定义小拟范畴构成的拟范畴 $\mathcal{Q}\mathrm{Cat}$，它建模小 $\infty$-范畴构成的 $\infty$-范畴 $\Cat_\infty$。

\section{构造的想法}

\nitem\label{L5-1.1} 上次课，对拟范畴 $\cC$ 与 $\cD$，我们构造了它们之间函子构成的拟范畴 $\Fun(\cC,\cD)$。设 $\Fun(\cC,\cD)^{\simeq}$ 为这个拟范畴的核心 Kan 复形，它可以看作从 $\Fun(\cC,\cD)$ 中丢弃全部不可逆自然变换而得到。

\nitem\label{L5-1.2} 利用这些 Kan 复形，我们可以把\textbf{普通}范畴 $\mathrm{QCat}$ 以 $\Set_\Delta$ 充实，得到一个单纯（充实）范畴 $\mathcal{Q}\mathrm{Cat}$。同伦假设提示：$\mathcal{Q}\mathrm{Cat}$ 已经是 $\Cat_\infty$ 的一个模型。

\nitem\label{L5-1.3} 为了把这件事翻译成严格的数学，我们将构造模型范畴之间的一个 Quillen 等价
\[
\mathfrak{C}\colon\Set^{\mathrm{Joyal}}_\Delta\rightleftarrows\Cat_\Delta:\mathfrak{N}_{\bullet},
\]
即单纯集的 Joyal 模型范畴与小单纯范畴的标准模型范畴之间的 Quillen 等价。单纯范畴 $\mathcal{Q}\mathrm{Cat}$ 是 $\Cat_\Delta$ 中的纤维对象\footnote{\linespread{1.45}\selectfont 严格说来，$\mathcal{Q}\mathrm{Cat}$ 不是 $\Cat_\Delta$ 中的对象，因为它不是小的。因此我们需要把小的单纯范畴构成的大模型范畴 $\Cat_\Delta$ 换成大的单纯范畴构成的非常大模型范畴 $\widetilde{\Cat}_\Delta$，并类似地把 $\Set^{\mathrm{Joyal}}_\Delta$ 换成大的单纯集构成的非常大的 Joyal 模型范畴 $\widetilde{\Set}^{\mathrm{Joyal}}_\Delta$。这里可以用 Grothendieck 宇宙给上述大小条件赋予确切含义。在真正遇到与之相关的问题之前，我们将忽略这些集合论问题。}，因为它以 $\Set_\Delta$ 中的纤维对象充实。Joyal 模型结构的设计使得（双）纤维对象恰好是拟范畴。因此
\[
\mathcal{Q}\mathrm{Cat}\defeq\mathfrak{N}_{\bullet}(\mathcal{Q}\mathrm{Cat})
\]
是 $\Set^{\mathrm{Joyal}}_\Delta$ 中的纤维对象，即一个拟范畴。

\section{拟范畴的同伦范畴}

\nitem\label{L5-2.1} 所希望的拟范畴 $\mathcal{Q}\mathrm{Cat}$ 的同伦范畴，按第 4 讲定义 4.2 的意义，可以容易地定义。它同时也会是所希望的模型范畴 $\Set^{\mathrm{Joyal}}_\Delta$ 的同伦范畴，按第 2 讲定义 2.19 的意义。它建模所希望的 $\infty$-范畴 $\Cat_\infty$ 的同伦范畴 $\mathrm{h}\Cat_\infty$。

\begin{definition}\label{L5-2.2}
设 $F_1,F_2\colon\cC\to\cD$ 是拟范畴之间的函子。我们把拟范畴 $\Fun(\cC,\cD)$ 中的态射 $\alpha\colon F_1\to F_2$ 称为从 $F_1$ 到 $F_2$ 的\textbf{自然变换}。

若 $\alpha$ 是 $\Fun(\cC,\cD)$ 中的同构，则称 $\alpha$ 是\textbf{可逆的}，或说 $\alpha$ 是一个\textbf{等价}。
\end{definition}

\begin{theorem}\label{L5-2.3}（Ker.01DK）
设 $F_1,F_2\colon\cC\to\cD$ 是拟范畴之间的函子。自然变换 $\alpha\colon F_1\to F_2$ 可逆，当且仅当它在任意对象 $x\in\cC$ 处的取值是同构，即 $\alpha(x)\colon F_1(x)\xrightarrow{\ \simeq\ }F_2(x)$。
\end{theorem}

\begin{exercise}\label{L5-2.4}
证明下述构造定义了一个普通范畴 $\mathrm{h}\mathcal{Q}\mathrm{Cat}$：
\begin{itemize}[leftmargin=2.2em,itemsep=0.1em]
\item $\mathrm{h}\mathcal{Q}\mathrm{Cat}$ 的对象是小拟范畴；
\item $\mathrm{h}\mathcal{Q}\mathrm{Cat}$ 的态射是拟范畴之间函子的等价类：
\[
\Hom_{\mathrm{h}\mathcal{Q}\mathrm{Cat}}(\cC,\cD)\defeq\pi_0(\Fun(\cC,\cD)^{\simeq});
\]
\item 复合由 $[G]\circ[F]\defeq[G\circ F]$ 给出。
\end{itemize}
\end{exercise}

\section{拟范畴之间的等价}

\nitem\label{L5-3.1} 注意我们有函子 $\mathcal{Q}\mathrm{Cat}\to\mathrm{h}\mathcal{Q}\mathrm{Cat}$。与第 4 讲定义 5.1 类似，我们作如下定义。

\begin{definition}\label{L5-3.2}
设 $F\colon\cC\to\cD$ 是拟范畴之间的函子，即 $\mathcal{Q}\mathrm{Cat}$ 中的态射。若 $F$ 在 $\mathrm{h}\mathcal{Q}\mathrm{Cat}$ 中的像 $[F]$ 是同构，则称 $F$ 是一个\textbf{等价}。
\end{definition}

\begin{remark}\label{L5-3.3}
换言之，$F\colon\cC\to\cD$ 是等价，当且仅当存在函子 $G\colon\cD\to\cX$，使得 $G\circ F$ 等价于 $\id_{\cC}$，且 $F\circ G$ 等价于 $\id_{\cD}$。
\end{remark}

\yizhu{原文此处的 $G\colon\cD\to\cX$ 中出现未加定义的记号 $\cX$，按上下文应为 $G\colon\cD\to\cC$；译文照录原文，未作改动。\textbf{[来源：I]}}

\begin{theorem}\label{L5-3.4}（Ker.01JX）
设 $F\colon\cC\to\cD$ 是拟范畴之间的函子。则 $F$ 是等价，当且仅当
\begin{enumerate}[label=(\arabic*),leftmargin=2.4em,itemsep=0.1em]
\item $F$ 是\textbf{全忠实的}：对任意对象 $x,y\in\cC$，函子 $\mathrm{Maps}_{\cC}(x,y)\to\mathrm{Maps}_{\cD}(Fx,Fy)$ 是 $\infty$-群胚之间的等价。
\item $F$ 是\textbf{本质满的}：映射 $\pi_0(\cC^{\simeq})\to\pi_0(\cD^{\simeq})$ 是满射。换言之，对任意对象 $d\in\cD$，存在 $c\in\cC$ 与同构 $F(c)\xrightarrow{\ \simeq\ }d$。
\end{enumerate}
\end{theorem}

\begin{exercise}\label{L5-3.5}
本质满性可以在同伦范畴的层面上检验。即：$F\colon\cC\to\cD$ 本质满，当且仅当 $\mathrm{h}F\colon\mathrm{h}\cC\to\mathrm{h}\cD$ 本质满。
\end{exercise}

\begin{exercise}\label{L5-3.6}
全忠实性不能在同伦范畴的层面上检验。
\end{exercise}

\begin{remark}\label{L5-3.7}
设 $F\colon\cC\to\cD$ 是拟范畴之间的函子。到目前为止，$F$ 上至少有四条与等价类似的条件：
\begin{enumerate}[label=(\alph*),leftmargin=2.4em,itemsep=0.1em]
\item $F$ 是 $\Set_\Delta$ 中的同构。
\item $F$ 是拟范畴之间的等价。
\item $F$ 是 Kan--Quillen 模型范畴 $\Set^{\mathrm{KQ}}_\Delta$ 中的弱同伦等价。
\item $F$ 诱导普通范畴之间的等价：$\mathrm{h}F\colon\mathrm{h}\cC\to\mathrm{h}\cD$。
\end{enumerate}
这些概念之间的关系是
\[
\begin{tikzcd}[column sep=3.4em,row sep=2.4em]
(a) \arrow[r, Rightarrow] & (b) \arrow[r, Rightarrow] \arrow[d, Rightarrow] & (c)\\
& (d) &
\end{tikzcd}
\]
\end{remark}

\begin{exercise}\label{L5-3.8}
上述蕴含关系都不可逆。
\begin{enumerate}[label=(\arabic*),leftmargin=2.4em,itemsep=0.1em]
\item $N_{\bullet}(\{\bullet\rightleftarrows\bullet\})\to\Delta^0$ 满足 (b) 但不满足 (a)。
\item $\Delta^1\to\Delta^0$ 满足 (c) 但不满足 (b) 与 (d)。
\item $\Sing(S^2)\to\Delta^0$ 满足 (d) 但不满足 (b) 与 (c)。
\end{enumerate}
\end{exercise}

\begin{remark}\label{L5-3.9}
设 $\mathrm{F}\colon\mathrm{C}\to\mathrm{D}$ 是由 $F\colon\cC\to\cD$ 建模的 $\infty$-范畴之间的函子。上述条件翻译为：
\begin{enumerate}[label=(\alph*),leftmargin=2.4em,itemsep=0.1em]
\item 这个条件没有意义，因为它是邪恶的。
\item $\mathrm{F}$ 是 $\infty$-范畴之间的等价。
\item 在把 $\mathrm{C}$ 与 $\mathrm{D}$ 中所有态射形式地变成可逆之后，$\mathrm{F}$ 成为 $\infty$-群胚之间的等价。
\item 在形式地迫使 $\mathrm{C}$ 与 $\mathrm{D}$ 中所有不可逆的高阶态射成为恒等之后，$\mathrm{F}$ 成为普通范畴之间的等价。
\end{enumerate}
\end{remark}

\begin{remark}\label{L5-3.10}
设 $\mathrm{C}$ 是 $\infty$-范畴。记 $\mathrm{C}^{\sharp}$ 为把 $\mathrm{C}$ 中所有态射形式地变成可逆而得到的 $\infty$-群胚。注意这不是核心 $\infty$-群胚 $\mathrm{C}^{\simeq}$——后者是丢弃所有不可逆态射而得到的（见第 4 讲练习 5.3）。我们有
\[
\mathrm{C}^{\simeq}\to\mathrm{C}\to\mathrm{C}^{\sharp}.
\]
\end{remark}

\section{Joyal 模型结构}

\nitem\label{L5-4.1} $\Set^{\mathrm{Joyal}}_\Delta$ 的类 (W) 中的态射将称为单纯集之间的\textbf{范畴等价}（categorical equivalence）\footnote{Joyal 称它们为弱范畴等价。我们这里沿用 Lurie 的术语。}。当限制到双纤维对象即拟范畴上时，它们应当建模底层的 $\infty$-范畴之间的等价。换言之：
\begin{center}
\emph{拟范畴之间的范畴等价，应当是定义~\ref{L5-3.2} 意义下的等价。}
\end{center}
因此主要任务是定义一般单纯集之间的范畴等价。

\nitem\label{L5-4.2} 回顾：在定义 $\Set_\Delta$ 上的经典模型结构时（见第 3 讲定理—定义 6.2），我们把单纯集之间的弱同伦等价宣告为这样的映射 $f\colon X\to Y$：$|f|\colon|X|\to|Y|$ 是 $\Top$ 中的弱同伦等价。用类似的方式定义 $\Set^{\mathrm{Joyal}}_\Delta$ 中的范畴等价是可能的。即：我们可以先定义 $\Cat_\Delta$ 中的弱等价，然后通过函子 $\mathfrak{C}\colon\Set_\Delta\to\Cat_\Delta$\footnote{这是 Lurie 的 HTT.2.2.5 中采用的做法，它与 Joyal 著作中的做法不同（见 [Joy08]）。更多信息见 HTT.2.2.5.9 与 HTT.2.2.5.10。}把它们搬运过来。然而，我觉得函子 $\mathfrak{C}$ 的构造不够直观，无法提供关于范畴等价的实际感受。因此我们将遵循 Joyal 的定义（见 [Joy08]）。

\nitem\label{L5-4.3} Joyal 的定义由下述观察所驱动：

\begin{exercise}\label{L5-4.4}
设 $\cC$ 是模型范畴，且其中每个对象都是上纤维的。则态射 $f\colon X\to Y$ 属于 (W)，当且仅当对任意纤维对象 $Z$，映射
\[
(\Hom_{\cC}(Y,Z)/{\sim})\to(\Hom_{\cC}(X,Z)/{\sim})
\]
是双射。这里“$\sim$”表示态射之间的同伦关系，见第 2 讲命题—定义 2.18。

提示：Yoneda 引理；函子 $\cC\to\cC[W^{-1}]$ 检测所有弱等价。
\end{exercise}

\begin{theoremdef}\label{L5-4.5}
设 $f\colon X\to Y$ 是 $\Set_\Delta$ 中的态射。下述条件等价：
\begin{enumerate}[label=(\arabic*),leftmargin=2.4em,itemsep=0.1em]
\item 对任意拟范畴 $Z$，映射
\[
\pi_0(\Fun(Y,Z)^{\simeq})\to\pi_0(\Fun(X,Z)^{\simeq})
\]
是集合之间的双射。
\item 对任意拟范畴 $Z$，函子
\[
\Fun(Y,Z)^{\simeq}\to\Fun(X,Z)^{\simeq}
\]
是 Kan 复形之间的弱等价。
\item 对任意拟范畴 $Z$，函子
\[
\Fun(Y,Z)\to\Fun(X,Z)
\]
是拟范畴之间的等价。
\end{enumerate}
若 $f$ 满足上述条件，则称 $f$ 是\textbf{范畴等价}。
\end{theoremdef}

\begin{proposition}\label{L5-4.6}（[Joy08, Corollary 2.29]）
任意内角嵌入 $\Lambda^n_i\to\Delta^n$（$0<i<n$）都是范畴等价。
\end{proposition}

\begin{exercise}\label{L5-4.7}
请说服自己：上述断言并不能由拟范畴的定义直接推出。
\end{exercise}

\begin{exercise}\label{L5-4.8}
范畴等价是弱同伦等价。

提示：第 3 讲 \S6.5。
\end{exercise}

\begin{exercise}\label{L5-4.9}（Ker.01EG）
无环 Kan 纤维化是范畴等价。
\end{exercise}

\begin{theoremdef}\label{L5-4.10}（[Joy08, Theorem 6.12]）
$\Set_\Delta$ 上存在一个模型结构，由以下各项给出：
\begin{itemize}[leftmargin=2.2em,itemsep=0.1em,label={}]
\item[(W)] 范畴等价构成的类；
\item[(C)] 单态射构成的类；
\item[(F)] 对 $(C\cap W)$ 具有右提升性质的态射构成的类。
\end{itemize}
我们称它为 $\Set_\Delta$ 上的\textbf{Joyal 模型结构}，并把这个模型范畴记作 $\Set^{\mathrm{Joyal}}_\Delta$。

这个模型范畴中的纤维对象恰好是拟范畴。
\end{theoremdef}

\nitem\label{L5-4.11} 下述结果是同义反复：

\begin{exercise}\label{L5-4.12}
函子 $\mathcal{Q}\mathrm{Cat}\to\Set_\Delta$ 诱导等价\footnote{事实上，它是同构。}$\mathrm{h}\mathcal{Q}\mathrm{Cat}\to\mathrm{h}\Set^{\mathrm{Joyal}}_\Delta$。特别地，$\Set^{\mathrm{Joyal}}_\Delta[W^{-1}]$ 也建模 $\mathrm{h}\Cat_\infty$。
\end{exercise}

\begin{definition}\label{L5-4.13}
$\Set^{\mathrm{Joyal}}_\Delta$ 中的纤维化称为\textbf{范畴纤维化}。
\end{definition}

\begin{proposition}\label{L5-4.14}
恒等函子给出一个 Quillen 伴随
\[
\Set^{\mathrm{Joyal}}_\Delta\rightleftarrows\Set^{\mathrm{KQ}}_\Delta.
\]
\end{proposition}

\begin{exercise}\label{L5-4.15}
设 $L\colon\Set^{\mathrm{Joyal}}_\Delta[W^{-1}]\rightleftarrows\Set^{\mathrm{KQ}}_\Delta[W^{-1}]:R$ 是上述 Quillen 伴随的导出伴随。把它解释为 $\mathrm{h}\Cat_\infty\rightleftarrows\mathrm{h}\Grpd_\infty$ 之间的一个伴随。
\end{exercise}

\begin{remark}\label{L5-4.16}
按定义，$\Set^{\mathrm{Joyal}}_\Delta$ 中的范畴纤维化 $p\colon X\to Y$ 对所有内角嵌入满足右提升性质。满足这些性质的态射称为\textbf{内纤维化}。注意当 $Y=\Delta^0$ 时，$Y$ 上的内纤维化就是范畴纤维化。然而，包含关系
\begin{center}
内纤维化 $\subset$ 范畴纤维化
\end{center}
一般而言是严格的。更多信息见 HTT.2.4.6.5。
\end{remark}

\section{单纯范畴}

\nitem\label{L5-5.1} 本节我们在小单纯范畴构成的范畴 $\Cat_\Delta$ 上定义一个模型结构。

\nitem\label{L5-5.2} 单纯范畴之间的弱等价，其定义与拓扑范畴之间的弱等价类似。见第 2 讲 \S3。

\begin{definition}\label{L5-5.3}
设 $\cC$ 是单纯范畴。它的\textbf{同伦范畴} $\pi_0\cC$ 由下式定义：
\[
\Ob(\pi_0\cC)\defeq\Ob(\cC),\qquad \Hom_{\pi_0\cC}(x,y)\defeq\pi_0(|\Hom_{\cC}(x,y)|).
\]
\end{definition}

\begin{definition}\label{L5-5.4}
设 $F\colon\cC\to\cD$ 是单纯范畴之间的函子。若满足以下两条，则称 $F$ 是\textbf{弱等价}：
\begin{itemize}[leftmargin=2.2em,itemsep=0.1em]
\item 它诱导等价 $\pi_0 F\colon\pi_0\cC\to\pi_0\cD$。
\item 对任意 $x,y\in\Ob(\cC)$，态射 $\Hom_{\cC}(x,y)\to\Hom_{\cD}(Fx,Fy)$ 是 $\Set^{\mathrm{KQ}}_\Delta$ 中的弱等价。
\end{itemize}
\end{definition}

\begin{theorem}\label{L5-5.5}（HTT.A.3.2）
$\Cat_\Delta$ 上存在一个典范模型结构，使得：
\begin{itemize}[leftmargin=2.2em,itemsep=0.1em]
\item 弱等价如定义~\ref{L5-5.4} 所述。
\item 纤维对象恰好是那些 $\Hom_{\cC}(-,-)$ 为 Kan 复形的单纯范畴 $\cC$。
\end{itemize}
我们称它为 $\Cat_\Delta$ 上的\textbf{经典}或\textbf{Bergner 模型结构}。
\end{theorem}

\begin{remark}\label{L5-5.6}
$\Cat_\Delta$ 中的纤维化与上纤维对象不容易描述，因此我们不打算描述它们。注意第 2 讲练习 3.6 有单纯模拟，它提示许多自然定义的单纯范畴不是上纤维的。
\end{remark}

\section{拟范畴与单纯范畴}

\nitem\label{L5-6.1} 本节我们构造所希望的 Quillen 等价
\[
\mathfrak{C}\colon\Set^{\mathrm{Joyal}}_\Delta\rightleftarrows\Cat_\Delta:\mathfrak{N}_{\bullet}.
\]
与 $|{-}|\colon\Set^{\mathrm{KQ}}_\Delta\rightleftarrows\Top:\Sing$ 的构造一样，我们只需沿 Yoneda 嵌入 $\Delta\to\Set_\Delta\defeq\Fun(\Delta^{\op},\Set)$ 找到限制 $\mathfrak{C}|_\Delta\colon\Delta\to\Cat_\Delta$ 的正确定义。整个函子 $\mathfrak{C}$ 将是唯一（在唯一等价意义下）扩张 $\mathfrak{C}|_\Delta$ 的余极限保持函子，而函子 $\mathfrak{N}_{\bullet}$ 将是它的右伴随。

\nitem\label{L5-6.2} 换言之，我们需要构造带有面映射与退化映射的单纯范畴 $\mathfrak{C}[\Delta^n]\in\Cat_\Delta$。为给这个构造提供动机，我们看下述问题：
\begin{center}
\emph{设 $\cC\in\Cat_\Delta$ 是纤维单纯范畴。函子 $\mathfrak{C}[\Delta^n]\to\cC$ 是什么？}
\end{center}
由于我们期望模型范畴 $\Cat_\Delta$ 建模 $\Cat_\infty$，又由于 $\mathfrak{C}[\Delta^n]$ 是上纤维的（因为我们期望 $\mathfrak{C}$ 是左 Quillen 函子），我们得到
\begin{center}
\emph{函子 $\mathfrak{C}[\Delta^n]\to\cC$ 应当建模相应 $\infty$-范畴之间的函子 $[n]\to\cC$。}
\end{center}

\nitem\label{L5-6.3} 例如，当 $n=2$ 时，它应当是见证 $\cC$ 中下述交换图的可逆 $2$-态射：
\[
\begin{tikzcd}[column sep=3.8em,row sep=2.4em]
& y \arrow[dl, "f"'] \arrow[dr, "g"] & \\
x \arrow[rr, "h"'] & & z.
\end{tikzcd}
\]
注意在单纯范畴 $\cC$ 中，复合是具体的：对 $0$-单形 $f\in\Hom_{\cC}(x,y)_0$ 与 $g\in\Hom_{\cC}(y,z)_0$，我们有良定义的单纯形 $g\circ f\in\Hom_{\cC}(x,z)_0$。因此上述可逆 $2$-态射对应于 $\Hom_{\cC}(x,z)_1$ 中的一个 $1$-单形 $g\circ f\to h$\footnote{也可以改用 $h\to g\circ f$。由于关于 $\Hom_{\cC}(x,z)$ 的纤维性假设，实际数据是一样的。但所得的函子 $\mathfrak{C}$ 会相差 $\Set_\Delta\to\Set_\Delta$、$X\mapsto X^{\op}$ 这个对合。这里我们遵循 Kerodon 中的约定（见 Ker.00KN），它与 HTT 中的约定相反（见 HTT.1.1.5.1）。}。

因此我们应当把 $\mathfrak{C}[\Delta^2]$ 定义为
\begin{itemize}[leftmargin=2.2em,itemsep=0.1em]
\item 它有三个对象，标号为 $0$、$1$ 与 $2$。
\item 态射单纯集为
\[
\Hom_{\mathfrak{C}[\Delta^2]}(i,j)\defeq
\begin{cases}
\varnothing, & \text{若 } i>j,\\
\Delta^1, & \text{若 } i=0,\ j=2,\\
\Delta^0, & \text{其他}.
\end{cases}
\]
\item 唯一不明显的复合是
\[
\Hom_{\mathfrak{C}[\Delta^2]}(0,1)\times\Hom_{\mathfrak{C}[\Delta^2]}(1,2)\to\Hom_{\mathfrak{C}[\Delta^2]}(0,2),
\]
即
\[
\Delta^0\times\Delta^0\to\Delta^1,
\]
它由 $0$-单形 $0\in\Delta^1$ 给出。
\end{itemize}

\nitem\label{L5-6.4} 上述研究可以对一般的 $n$ 进行。即：对任意 $\infty$-范畴，由于 $[n]$ 来自一个线性序集\footnote{另一方面，第 1 讲第 3 节中讨论的同伦连贯交换图超出了我们的能力范围，因为它们不是由（偏）序集给出的。}，我们知道函子 $[n]\to\mathrm{C}$ 的正确含义。例如，$\Hom_{\mathfrak{C}[\Delta^n]}(0,n)$ 中的 $0$-单形应当是包含 $0$ 与 $n$ 的子集 $I\subset[n]$；而每个链 $I_0\supset\cdots\supset I_m$ 对应于一个 $m$-单形。

用公式写出，我们定义：

\begin{definition}\label{L5-6.5}
设 $J$ 是任意线性序集。我们如下定义一个单纯范畴 $\mathfrak{C}[\Delta^J]$：
\begin{itemize}[leftmargin=2.2em,itemsep=0.1em]
\item 对象是 $J$ 中的元素；
\item 态射单纯集为
\[
\Hom_{\mathfrak{C}[\Delta^J]}(x,y)\defeq
\begin{cases}
\varnothing, & \text{若 } x>y,\\
\mathfrak{N}_{\bullet}(P_{x,y}), & \text{若 } x\le y,
\end{cases}
\]
其中 $P_{x,y}\defeq\{I\subset[x,y]\mid x,y\in I\}$ 配备偏序 $\supset$，而 $\mathfrak{N}_{\bullet}(P_{x,y})$ 是相应范畴的脉；
\item 态射按下式诱导：
\[
(I\subset[x,y],\ I'\subset[y,z])\mapsto(I\cup I'\subset[x,z]).
\]
\end{itemize}
\end{definition}

\nitem\label{L5-6.6} 上述构造对 $J$ 是函子性的，因此我们有函子
\[
\Delta\to\Cat_\Delta,\qquad [n]\mapsto\mathfrak{C}[\Delta^n]\defeq\mathfrak{C}[\Delta^{[n]}].
\]
注意 $\mathfrak{C}[\Delta^n]$ 不是纤维的。这一点在 $n=2$ 时已经可以看出。设
\[
\mathfrak{C}\colon\Set_\Delta\to\Cat_\Delta
\]
是本质唯一地扩张上述函子的余极限保持函子。

\begin{exercise}\label{L5-6.7}
$\mathfrak{C}[\Lambda^2_1]$ 是什么？
\end{exercise}

\begin{theoremdef}\label{L5-6.8}（HTT.2.2.5.1, HTT.2.2.5.8）
函子 $\mathfrak{C}\colon\Set_\Delta\to\Cat_\Delta$ 有右伴随，称为\textbf{单纯脉函子}
\[
\mathfrak{N}_{\bullet}\colon\Cat_\Delta\to\Set_\Delta.
\]
伴随对
\[
\mathfrak{C}\colon\Set^{\mathrm{Joyal}}_\Delta\rightleftarrows\Cat_\Delta:\mathfrak{N}_{\bullet}
\]
是 Quillen 等价。
\end{theoremdef}

\yizhu{原书在不同处混用罗马体 $N_{\bullet}$ 与花体 $\mathfrak{N}_{\bullet}$，二者同指单纯脉函子（见定理—定义~\ref{L5-6.8}）；译文照录原文的字体区分。\textbf{[来源：I]}}

\nitem\label{L5-6.9} 注意 $\mathfrak{C}[\Delta^0]\simeq\{\ast\}$。因此我们可以
\begin{itemize}[leftmargin=2.2em,itemsep=0.1em]
\item 对任意单纯集 $X$，把 $X_0$ 与 $\Ob(\mathfrak{C}(X))$ 等同；
\item 对任意单纯范畴 $\cC$，把 $\Ob(\cC)$ 与 $\mathfrak{N}_{\bullet}(\cC)_0$ 等同。
\end{itemize}
从现在起，我们将滥用记号，把它们看作同一个集合。

\nitem\label{L5-6.10} 为了说明 $\Set^{\mathrm{Joyal}}_\Delta$ 与 $\Cat_\Delta$ 中相应的双纤维对象确实建模同一个 $\infty$-范畴，我们需要下述结果。

\begin{theorem}\label{L5-6.11}（Ker.01LE）
设 $\cC\in\Cat_\Delta$ 是纤维单纯范畴。则 Kan 复形之间存在典范弱同伦等价
\[
\Hom_{\cC}(x,y)\xrightarrow{\ \simeq\ }\Hom^{L}_{\mathfrak{N}_{\bullet}(\cC)}(x,y).
\]
\end{theorem}

\begin{corollary}\label{L5-6.12}
设 $\cC\in\Cat_\Delta$ 是纤维单纯范畴。则存在典范等价 $\mathrm{h}\cC\xrightarrow{\ \simeq\ }\mathrm{h}\mathfrak{N}_{\bullet}(\cC)$\footnote{事实上，这个结果要初等得多。见 Ker.00M4。}。
\end{corollary}

\section{拟范畴的拟范畴}

\nitem\label{L5-7.1} 本节的文字是 \S1 的逐字重复。

\begin{definition}\label{L5-7.2}
设 $\mathcal{Q}\mathrm{Cat}$ 是由以下各项定义的单纯范畴：
\begin{itemize}[leftmargin=2.2em,itemsep=0.1em]
\item 对象是拟范畴；
\item 对拟范畴 $\cC$ 与 $\cD$，
\[
\Hom_{\mathcal{Q}\mathrm{Cat}}(\cC,\cD)\defeq\Fun(\cC,\cD)^{\simeq};
\]
\item 复合
\[
\Fun(\cC,\cD)\times\Fun(\cD,\cE)\to\Fun(\cC,\cE)
\]
由 $\Fun(-,-)$ 的泛性质给出。
\end{itemize}
\end{definition}

\begin{definition}\label{L5-7.3}
\textbf{小拟范畴构成的拟范畴}定义为
\[
\mathcal{Q}\mathrm{Cat}\defeq\mathfrak{N}_{\bullet}(\mathcal{Q}\mathrm{Cat}).
\]
设 $\Cat_\infty$ 是由它建模的 $\infty$-范畴\footnote{用严格的数学语言，我们应当把 $\infty$-范畴定义为拟范畴。然而，我确实希望读者把拟范畴视为 $\infty$-范畴的模型。}。我们称它为\textbf{小 $\infty$-范畴构成的 $\infty$-范畴}。
\end{definition}

\yizhu{原书用两种字体区分两个不同的对象：黑板体的 $\mathcal{Q}\mathrm{Cat}$ 是小拟范畴构成的（单纯）范畴，花体的 $\mathcal{Q}\mathrm{Cat}$ 是由它得到的拟范畴（见定义~\ref{L5-7.2}、\ref{L5-7.3}）。附录标题“QCat vs. QCat”正指这两者，译文沿用这一区分。\textbf{[来源：I]}}

\begin{exercise}\label{L5-7.4}
设 $\cC$ 与 $\cD$ 是拟范畴。Kan 复形之间存在典范弱同伦等价
\[
\Fun(\cC,\cD)^{\simeq}\xrightarrow{\ \simeq\ }\Hom^{L}_{\mathcal{Q}\mathrm{Cat}}(\cC,\cD).
\]
换言之，$\Fun(\cC,\cD)^{\simeq}$ 如所愿地建模 $\infty$-\textbf{群胚} $\mathrm{Maps}_{\Cat_\infty}(\cC,\cD)$。
\end{exercise}

\begin{variant}\label{L5-7.5}
类似地，我们定义\textbf{小 Kan 复形构成的拟范畴}
\[
\mathcal{K}\mathrm{an}\defeq\mathfrak{N}_{\bullet}(\mathbb{K}\mathrm{an}),
\]
其中 $\mathbb{K}\mathrm{an}\subset\mathcal{Q}\mathrm{Cat}$ 是由 Kan 复形构成的全单纯子范畴。与之前一样，拟范畴 $\mathcal{K}\mathrm{an}$ 建模小 $\infty$-群胚构成的 $\infty$-范畴 $\Grpd_\infty$。
\end{variant}

\begin{exercise}\label{L5-7.6}
证明 $\mathcal{Q}\mathrm{Cat}$ 的同伦范畴确实等价于练习~\ref{L5-2.4} 中的 $\mathrm{h}\mathcal{Q}\mathrm{Cat}$。
\end{exercise}

\section{单纯范畴中的交换图}

\nitem\label{L5-8.1} 由前面的讨论，在 $\infty$-范畴论中，对任意单纯集 $S$，以 $S$ 为指标的交换图在单纯范畴 $\cC$ 中的正确概念，是函子 $\mathfrak{C}[S]\to\cC$。我们把这样的图称为 $\cC$ 中以 $S$ 为指标的\textbf{同伦连贯交换图}。

\nitem\label{L5-8.2} 在大多数情形下，图的形状是一个普通范畴 $\cD$，那时我们也滥用语言，把函子 $\mathfrak{C}[N_{\bullet}(\cD)]\to\cC$ 称为 $\cC$ 中以 $\cD$ 为指标的同伦连贯交换图。这个概念在标准的 $\infty$-范畴构造下表现良好。

人们也可以要求图\textbf{严格交换}，即要求一个函子 $\cD\to\cC$，其中 $\cD$ 被看作以离散单纯集充实。在实践中，构造这样的图是困难的，但大多数非形式计算最终会把人们引向这一领域\footnote{我的个人经验：真正的非形式计算非常少。}。

人们也可以要求图\textbf{同伦交换}，即要求一个函子 $\cD\to\mathrm{h}\cC$。这本质上是普通范畴论中的概念。

\nitem\label{L5-8.3} 我们有：
\begin{center}
严格交换 $\Rightarrow$ 同伦连贯交换 $\Rightarrow$ 同伦交换。
\end{center}

\begin{exercise}\label{L5-8.4}
把上述讨论应用于 $\cC\defeq\mathcal{Q}\mathrm{Cat}$，得到拟范畴的交换图的三种概念。
\end{exercise}

\section{什么是 $\infty$-范畴？}

\nitem\label{L5-9.1} 在这些讲义中，我们采取下述观点：
\begin{itemize}[leftmargin=2.2em,itemsep=0.1em]
\item 一个 $\infty$-范畴由一个拟范畴建模，而两个拟范畴之间的任何等价都给出它们所建模的 $\infty$-范畴之间的一个等同。
\end{itemize}
然而注意，这样的等同并不唯一，而且等同之间还可以有等同，等等。

特别地，我们不把某个特定的拟范畴指派给一个 $\infty$-范畴\footnote{然而注意，在 Lurie 的书中，$\infty$-范畴就是拟范畴。}。例如，谈论 $\infty$-范畴中对象的集合是没有意义的。相反，人们可以谈论对象的 $\infty$-群胚（因为 $\infty$-群胚是在同样的哲学下定义的\footnote{见第 4 讲注 8.4。}），从而谈论对象的等价类的集合。

\nitem\label{L5-9.2} 为了用更严格的语言描述上述哲学，我们作如下定义。

\begin{definition}\label{L5-9.3}
设 $\cC$ 是拟范畴，$W$ 是 $\cC$ 中一族态射。若对任意拟范畴 $\cE$，函子
\[
\Fun(\cD,\cE)\to\Fun(\cC,\cE)
\]
是全忠实的，且其本质像恰好包含把 $W$ 中的态射送到同构的函子 $\cC\to\cE$，则称函子 $\cC\to\cD$ 把 $\cD$ 实现为 $\cC$ 关于 $W$ 的\textbf{拟范畴局部化}。
\end{definition}

\nitem\label{L5-9.4} 可以证明拟范畴局部化本质唯一，即在可缩的选择空间意义下唯一（见 HTT.3.1）。我们将用 $\cC[W^{-1}]$ 表示 $\cC$ 关于 $W$ 的拟范畴局部化。由定理~\ref{L5-A.3}，我们有典范等价
\[
\mathfrak{N}_{\bullet}(\mathcal{Q}\mathrm{Cat})[W^{-1}]\xrightarrow{\ \simeq\ }\mathcal{Q}\mathrm{Cat},
\]
其中 $W$ 是范畴等价构成的族。

\nitem\label{L5-9.5} 现在我们作如下定义：
\begin{itemize}[leftmargin=2.2em,itemsep=0.1em]
\item 一个 $\infty$-范畴是 $\mathfrak{N}_{\bullet}(\mathcal{Q}\mathrm{Cat})[W^{-1}]$ 关于范畴等价的拟范畴局部化中的一个对象。
\item 对拟范畴 $\cC\in\mathcal{Q}\mathrm{Cat}$，由 $\cC$ 建模的 $\infty$-范畴是它在 $\mathfrak{N}_{\bullet}(\mathcal{Q}\mathrm{Cat})\to\mathfrak{N}_{\bullet}(\mathcal{Q}\mathrm{Cat})[W^{-1}]$ 下的像。
\end{itemize}
在上文中，我们不对 $\mathfrak{N}_{\bullet}(\mathcal{Q}\mathrm{Cat})[W^{-1}]$ 作具体选择，而本讲义中关于 $\infty$-范畴的任何陈述都不会依赖于这样的选择。特别地，我们不把 $\infty$-范畴看作 $\mathcal{Q}\mathrm{Cat}$ 中的 $0$-单形。

模去脚注 1 中的问题，$\mathfrak{N}_{\bullet}(\mathcal{Q}\mathrm{Cat})[W^{-1}]$ 本身建模一个 $\infty$-范畴，它称为\textbf{（小）$\infty$-范畴构成的 $\infty$-范畴}，记作 $\Cat_\infty$。

\nitem\label{L5-9.6} 在实践中，上述哲学导致下述原则。要作关于 $\infty$-范畴的定义或构造，我们可以：
\begin{enumerate}[label=(\roman*),leftmargin=2.4em,itemsep=0.1em]
\item 通过拟范畴来作，并证明所得结果在范畴等价下不变。
\item 组合已有的关于 $\infty$-范畴的定义或构造，而不提及拟范畴。
\end{enumerate}
当两种方法都可用时，我们偏爱第二种，因为它允许我们把定义与构造翻译到 $\infty$-范畴论的其他进路中。

\appendixsection{$\mathcal{Q}\mathrm{Cat}$ 与 $\mathcal{Q}\mathrm{Cat}$}

\begin{exercise}\label{L5-A.1}
设 $\cC$ 是普通范畴。构造同构 $N_{\bullet}(\cC)\cong\mathfrak{N}_{\bullet}(\cC)$，其中右端把 $\cC$ 看作以离散单纯集充实。
\end{exercise}

\begin{exercise}\label{L5-A.2}
构造函子 $\mathfrak{N}_{\bullet}(\mathcal{Q}\mathrm{Cat})\to\mathcal{Q}\mathrm{Cat}$，使其建模 $\mathcal{Q}\mathrm{Cat}\to\Cat_\infty$。
\end{exercise}

\begin{theorem}\label{L5-A.3}
函子 $\mathfrak{N}_{\bullet}(\mathcal{Q}\mathrm{Cat})\to\mathcal{Q}\mathrm{Cat}$ 把 $\mathcal{Q}\mathrm{Cat}$ 实现为 $\mathfrak{N}_{\bullet}(\mathcal{Q}\mathrm{Cat})$ 关于范畴等价类的拟范畴局部化。
\end{theorem}

\nitem\label{L5-A.4} \textbf{建议阅读。} [DK80c]、[DK80a]、[DK80b]（原始文献）与 [Hin16]（较新）。

\subsection*{参考文献}
\par\noindent\hangindent=5.2em\hangafter=1 [DK80a]\hspace{0.6em}William G Dwyer and Daniel M Kan. Calculating simplicial localizations. \emph{Journal of Pure and Applied Algebra}, 18(1):17--35, 1980.
\par\noindent\hangindent=5.2em\hangafter=1 [DK80b]\hspace{0.6em}William G Dwyer and Daniel M Kan. Function complexes in homotopical algebra. \emph{Topology}, 19(4):427--440, 1980.
\par\noindent\hangindent=5.2em\hangafter=1 [DK80c]\hspace{0.6em}William G Dwyer and Daniel M Kan. Simplicial localizations of categories. \emph{Journal of pure and applied algebra}, 17(3):267--284, 1980.
\par\noindent\hangindent=5.2em\hangafter=1 [Hin16]\hspace{0.6em}Vladimir Hinich. Dwyer--kan localization revisited. \emph{Homology, Homotopy and Applications}, 18(1):27--48, 2016.
\par\noindent\hangindent=5.2em\hangafter=1 [Joy08]\hspace{0.6em}Andr\'e Joyal. The theory of quasi-categories and its applications. 2008.
